\chapter{Configurations et preuves géométriques}
\section*{Exercice 11.1 -- Milieux}
$\overrightarrow{MN}=\overrightarrow{AN}-\overrightarrow{AM}=\tfrac12(\overrightarrow{AC}-\overrightarrow{AB})=\tfrac12\overrightarrow{BC}$. On obtient le parallélisme et le rapport de longueurs $1/2$.

\section*{Exercice 11.2 -- Barycentre}
La somme des poids vaut $3$. Les coordonnées pondérées donnent $G((0+12)/3,(2-2)/3)=(4,0)$. Vérification : $\overrightarrow{GA}=(-4,2)$ et $2\overrightarrow{GB}=(4,-2)$ somment au vecteur nul.

\section*{Exercice 11.3 -- Médianes}
L'isobarycentre est $G(2,3)$. Le milieu de $[BC]$ est $I(3,4{,}5)$ ; $\overrightarrow{AG}=(2/3)\overrightarrow{AI}$, donc $G$ est aux deux tiers depuis $A$.

\section*{Exercice 11.4 -- Cercle circonscrit}
Le triangle est rectangle en $A$. Les médiatrices des côtés sur les axes sont $x=2$ et $y=1{,}5$, donc le centre est $(2,1{,}5)$. Sa distance à $A$ vaut $\sqrt{4+2{,}25}=2{,}5$, moitié de l'hypoténuse $5$.

\section*{Exercice 11.5 -- Aire}
La hauteur associée au côté $5$ vaut $8\sin30\degre=4$. L'aire est donc $(5\times4)/2=10$. C'est $\frac12 ab\sin\theta$, obtenue à partir de base fois hauteur divisée par deux.

\section*{Exercice 11.6 -- Lieu}
L'égalité de distances équivaut à $x^2+(y-2)^2=y^2$, donc $y=x^2/4+1$. Pour $(2,2)$, la distance à $F$ vaut $2$, comme celle à l'axe ; il appartient bien à la parabole.
